\section{Workflow with fine-tuning of transformer models}
\label{sec:finetuning}

Fine-tuning of pretrained transformer models can be important for the application at hand.
In the proposed flow of entity resolution, this concerns the bi-encoder model for embedding
and the cross-encoder model for direct text-to-link classification.


\subsection{Fine-tuning transformer models}
\label{sec:fine-tuning-models}

Using the training-test construction described in Section
\ref{sec:training-candidates}, we now consider task-specific fine-tuning
of the bi-encoder and cross-encoder models. The same random splits and
treatment of multiple valid partners are retained, so that the results
can be compared directly with those obtained using the pretrained models.

There is one important distinction between the two fine-tuning
experiments. After fine-tuning the bi-encoder, the records are re-embedded
and the $k=50$ candidate blocks are reconstructed from the fine-tuned
embeddings before the traditional scoring models are evaluated. For the
cross-encoder, the candidate blocks obtained from the pretrained
bi-encoder are held fixed. Thus, the cross-encoder comparison isolates
the effect of fine-tuning the classifier rather than changing the
candidate set at the same time.

The hyperparameters for learning transformer models involve at least four choices: the
warmup step, typically set to around 10\%; the learning rate, typically of the order
$10^{-5}$; the number of epochs, typically small; and the batch size, typically 16 or 32.
While some experimentation is needed to settle the hyperparameters in practice, it is rare
to apply cross-validation inside the training set formally due to the computation cost this
would require. Notice that the computation cost differs for bi-encoder and cross-encoder.
Once the embeddings from a bi-encoder have been obtained, similarity calculation and
subsequent scoring are computationally cheap. In contrast, the cross-encoder must process
each pair of texts directly whenever it is applied.

The bi-encoder model \texttt{all-MiniLM-L6-v2} is trained for one epoch with
\texttt{ContrastiveLoss}, using 10\% linear warmup, batch size 16 and learning rate
$2\times10^{-5}$. The contrastive loss encourages the embeddings of matched record
pairs to be closer and those of non-matched pairs to be further apart. The fine-tuned
bi-encoder model can then be applied to the records in $A\cup B$ to yield the fine-tuned
embeddings, from which the per-field cosine similarities are calculated as before.

The cross-encoder \texttt{ms-marco-MiniLM-L-6-v2} is trained for five epochs with
10\% linear warmup, batch size 16, learning rate $2\times10^{-5}$, and a maximum
sequence length of 128 tokens. The training pairs are labelled as matches or non-matches
for binary link classification.

For comparison, we also evaluate the pretrained cross-encoder without updating any of
its model parameters. The same 50--50 outer split and 80--20 pair-level validation split
are used, but the training portion is not used to fit the cross-encoder. It is retained only
so that the link-classification threshold can be selected on exactly the same validation
partition as for the fine-tuned model. In both cases, the threshold is selected from
$0.10,0.15,\ldots,0.90$ by maximising the validation F1 score, and the direct output of
\texttt{CrossEncoder.predict()} is used without an additional sigmoid transformation.

While the hyperparameters above may not be optimal for these models, they suffice for
the methodological investigations below, which is our purpose in this paper rather than
to find the best procedures for the two datasets at hand.

\rev{The bi-encoder and cross-encoder are fine-tuned for one and five epochs,
respectively. Thus, the fine-tuned comparisons in Table~\ref{tab:fine-tuning}
concern the specified training configurations, rather than an epoch-matched
comparison of model architectures. The contribution of training duration to
the observed differences is not assessed separately.}

\begin{table}[ht]
\centering
\caption{FLR and MMR for pretrained and fine-tuned transformer models. For the
generative and logistic models, pretrained and fine-tuned refer to the
bi-encoder used for embedding and blocking; the candidate blocks are
reconstructed after bi-encoder fine-tuning. For the cross-encoder, they refer
to the cross-encoder model itself, with the candidate blocks held fixed.
\rev{Entries give the mean (Monte Carlo standard error) over three random splits.}}
\scriptsize
\setlength{\tabcolsep}{3pt}
\begin{tabular}{llrrrr}
\toprule
& & \multicolumn{2}{c}{Pretrained} & \multicolumn{2}{c}{Fine-tuned} \\
\cmidrule(lr){3-4}\cmidrule(lr){5-6}
Dataset & Model & FLR & MMR & FLR & MMR \\
\midrule
DBLP-Scholar & Generative, kernel density & 0.043 \rev{(0.0026)} & 0.014 \rev{(0.0009)} & 0.046 \rev{(0.0044)} & 0.011 \rev{(0.0012)} \\
 & Discriminative, logistic & 0.009 \rev{(0.0027)} & 0.068 \rev{(0.0027)} & 0.026 \rev{(0.0040)} & 0.027 \rev{(0.0042)} \\
 & Cross-encoder & 0.061 \rev{(0.0026)} & 0.026 \rev{(0.0011)} & 0.018 \rev{(0.0031)} & 0.014 \rev{(0.0009)} \\
\addlinespace
Abt-Buy & Generative, kernel density & 0.327 \rev{(0.0043)} & 0.356 \rev{(0.0073)} & 0.124 \rev{(0.0008)} & 0.145 \rev{(0.0031)} \\
 & Discriminative, logistic & 0.233 \rev{(0.0210)} & 0.543 \rev{(0.0597)} & 0.053 \rev{(0.0023)} & 0.269 \rev{(0.0059)} \\
 & Cross-encoder & \rev{0.267} \rev{(0.0102)} & \rev{0.278} \rev{(0.0123)} & 0.041 \rev{(0.0023)} & 0.152 \rev{(0.0130)} \\
\bottomrule
\end{tabular}
\label{tab:fine-tuning}
\end{table}

Table \ref{tab:fine-tuning} compares the link classification performance with or without
fine-tuning of the transformer models. The results for the traditional models with
pretrained bi-encoder embeddings are reproduced from Table \ref{tab:comparison} for easy
comparison. All results are averaged over three random 50--50 splits of the training and
test sets.

The cross-encoder results show a clear improvement from fine-tuning. For DBLP-Scholar,
the pretrained cross-encoder has an FLR of 0.061 and an MMR of 0.026, compared with
0.018 and 0.014 after fine-tuning. For Abt-Buy, the corresponding error rates fall from
\rev{0.267} and \rev{0.278} to 0.041 and 0.152. Fine-tuning therefore reduces both error rates for
the cross-encoder in both datasets.

Compared with the traditional models based on the fine-tuned bi-encoder, the
\rev{cross-encoder has a lower observed FLR than MEC in both datasets. For DBLP-Scholar,
the mean MMRs for MEC and the cross-encoder are 0.011 and 0.014, respectively;
this small difference should be interpreted cautiously given the Monte Carlo
standard errors. For Abt-Buy, their mean MMRs are likewise close (0.145 and
0.152), while their FLRs differ more substantially (0.124 and 0.041).
Logistic regression has higher observed MMR in both datasets, while its FLR is
closer to that of the cross-encoder for Abt-Buy.}

We also examined the sensitivity of the fine-tuned cross-encoder to the
block size. For Abt-Buy, reducing the block size from $k=50$ to $k=5$ increases
both the FLR and MMR, from 0.041 and 0.152 to 0.056 and 0.219, respectively.
For DBLP-Scholar, the effect is less uniform: the FLR increases from 0.018 to
0.022, whereas the MMR changes only slightly, from 0.014 to 0.013.

Changing $k$ affects more than the number of non-match examples available during
fine-tuning. It also changes the candidate set available at evaluation and hence the
probability that at least one valid partner is retrieved. This is particularly relevant
for Abt-Buy, where the floor MMR is appreciably higher at $k=5$ than at $k=50$.
The comparison should therefore be interpreted as sensitivity to the block-size
configuration as a whole, rather than as an isolated effect of the number of training
non-matches.

This points to another aspect of the cost-accuracy tradeoff for
cross-encoder models. Larger candidate blocks provide more retrieved hard
non-matches for supervised learning and can increase the chance that a valid
partner is available for classification, but they also increase the number of
record pairs to which the cross-encoder must be applied. The appropriate block
size therefore depends on candidate recall, linkage error and computational cost
together. Notice also that the test-set performance of a fine-tuned cross-encoder
would not necessarily carry over unchanged if a different blocking procedure were
employed for out-of-sample records.

Meanwhile, the workflow using the fine-tuned bi-encoder produces a different
pattern across the two datasets. Since the candidate blocks are reconstructed after
fine-tuning, this comparison reflects both the changed embedding representation and
the candidate blocks obtained from it. For DBLP-Scholar, the workflow using the
fine-tuned bi-encoder yields lower MMRs for both the generative and discriminative
models, but higher FLRs. For Abt-Buy, it yields considerably lower values of both
error rates for both traditional models. \rev{For the discriminative logistic regression model, the observed Abt-Buy FLR is
0.053, compared with 0.041 for the fine-tuned cross-encoder. This motivates
examining whether traditional-model scoring can handle some cases without
cross-encoder evaluation, which we investigate next.}

\subsection{Combining models for scoring-classification}
\label{sec:combining-models}

Since bi-encoder blocking does not distinguish different blocks of text pairs, which are
all sent to the cross-encoder model for link classification, a natural idea is to use the
traditional models to separate the harder and easier cases, each composed of $(i,C_i)$,
so that only the harder cases are passed on to the costly cross-encoder link classification,
while the easier cases can be classified by the traditional models directly.

Given a traditional discriminative model for scoring in Equation
\eqref{eq:match-probability}, trained on the fine-tuned bi-encoder embeddings, define a
corresponding discrimination gap for each case $(i,C_i)$, which is given as

\begin{equation}
\label{eq:discrimination-gap}
\Delta_i = p_{(i1)} - p_{(i2)},
\end{equation}

where $p_{(i1)} \geq p_{(i2)} \geq \cdots \geq p_{(ik)}$ are the order statistics of
$\{p_{ij}:j\in C_i\}$. A smaller gap \eqref{eq:discrimination-gap} for a file-$A$
record suggests lower confidence by the traditional model to separate its best file-$B$
candidate from the runner-up, in which sense link classification of this case is harder
than another case with a larger discrimination gap.

For the comparison below, the $k=50$ candidate blocks are held fixed at those obtained
from the pretrained bi-encoder, so that the logistic regression and cross-encoder are
evaluated on the same candidate records. The logistic regression scores are recalculated
using the fine-tuned bi-encoder embeddings. The quartiles are then formed separately
within each test set according to $\Delta_i$. \rev{These test-set quartiles are used only
as an illustrative analysis of how linkage errors vary with case difficulty; they do not
define a routing threshold for deployment. In an operational combined procedure, the
routing rule, including any discrimination-gap threshold, would be selected using the
training or validation partition and then applied unchanged to test or out-of-sample
records.}

\begin{table}[ht]
\centering
\caption{FLR and MMR by discrimination-gap quartiles and models. \rev{Entries give
the mean (Monte Carlo standard error) over three random 50--50 training-test
splits. A zero MCSE denotes no observed variation across these splits.}}
\scriptsize
\setlength{\tabcolsep}{3pt}
\begin{tabular}{llrrrr}
\toprule
& & \multicolumn{2}{c}{Logistic regression}
& \multicolumn{2}{c}{Cross-encoder} \\
\cmidrule(lr){3-4}\cmidrule(lr){5-6}
Dataset & Quartile & FLR & MMR & FLR & MMR \\
\midrule
DBLP-Scholar & Q1 (hardest) & 0.037 \rev{(0.0150)} & 0.045 \rev{(0.0042)} & 0.033 \rev{(0.0027)} & 0.029 \rev{(0.0010)} \\
 & Q2 & 0.031 \rev{(0.0062)} & 0.063 \rev{(0.0026)} & 0.034 \rev{(0.0100)} & 0.022 \rev{(0.0010)} \\
 & Q3 & 0.006 \rev{(0.0019)} & 0.041 \rev{(0.0080)} & 0.008 \rev{(0.0010)} & 0.010 \rev{(0.0027)} \\
 & Q4 (easiest) & 0.000 \rev{(0.0000)} & 0.000 \rev{(0.0000)} & 0.000 \rev{(0.0000)} & 0.000 \rev{(0.0000)} \\
\addlinespace
Abt-Buy & Q1 (hardest) & 0.230 \rev{(0.0229)} & 0.706 \rev{(0.0153)} & 0.121 \rev{(0.0085)} & 0.400 \rev{(0.0343)} \\
 & Q2 & 0.101 \rev{(0.0031)} & 0.259 \rev{(0.0559)} & 0.051 \rev{(0.0114)} & 0.131 \rev{(0.0150)} \\
 & Q3 & 0.017 \rev{(0.0025)} & 0.017 \rev{(0.0025)} & 0.016 \rev{(0.0045)} & 0.072 \rev{(0.0108)} \\
 & Q4 (easiest) & 0.000 \rev{(0.0000)} & 0.000 \rev{(0.0000)} & 0.000 \rev{(0.0000)} & 0.005 \rev{(0.0025)} \\
\bottomrule
\end{tabular}
\label{tab:discrimination-gap}
\end{table}

Table \ref{tab:discrimination-gap} shows the FLR and MMR by the quartiles of
$\Delta_i$, where Q1 refers to the smallest 25\% discrimination gaps and the hardest
cases accordingly, and Q4 refers to the 25\% easiest cases. The FLR and MMR of the
fine-tuned cross-encoder model are compared to those of logistic regression based on the
fine-tuned bi-encoder embeddings.

The results support the discrimination gap as a useful measure for separating harder
cases from easier ones, although the error rates do not decrease monotonically across
every quartile. The distinction is particularly clear for Abt-Buy. In Q1, logistic
regression has FLR 0.230 and MMR 0.706, compared with 0.121 and 0.400,
respectively, for the cross-encoder. The advantage of the cross-encoder remains
substantial in Q2. In Q3, however, both models have very low FLRs and logistic
regression has the lower MMR. In Q4, logistic regression makes no errors, while the
cross-encoder has only a very small MMR.

\rev{For DBLP-Scholar, neither model makes an observed error in Q4 across these splits.
In Q1--Q3, the cross-encoder has lower mean MMR, while the small FLR differences
in Q1--Q3 should not be treated as stable model orderings given the Monte Carlo
standard errors. The results are consistent with concentrating cross-encoder
evaluation among cases with smaller discrimination gaps, while leaving many of
the easier cases to the traditional model.}

This shows that, (i) the discrimination gap is a reasonable measure for separating the
harder cases from the easier ones, and (ii) it can be more beneficial to apply the
cross-encoder model to the harder cases in order to improve accuracy subject to a cost
constraint, provided the traditional models can do a relatively good job for the easier
cases. The Abt-Buy results illustrate this particularly clearly: the largest gains from the
cross-encoder occur among the cases with the smallest discrimination gaps, whereas
there is little benefit from sending the easiest cases to the more costly model.

\rev{Figure \ref{fig:flr-mmr-summary} provides a visual summary of the FLR/MMR
tradeoff in the main comparisons and the discrimination-gap analysis. The figure is
descriptive and intended to aid visual comparison; the corresponding tables remain
the primary source for the exact numerical values and Monte Carlo standard errors.}

\begin{figure}[ht]
\centering
\includegraphics[width=0.98\textwidth]{flr_mmr_tradeoff.pdf}
\caption{\rev{Visual summary of the false linkage rate (FLR) and missing match rate
(MMR) tradeoff. Panels (a) and (b) summarise the main comparisons reported in
Tables~\ref{tab:comparison}--\ref{tab:fine-tuning} for DBLP-Scholar and Abt-Buy,
respectively. Panels (c) and (d) summarise the discrimination-gap quartile results
from Table~\ref{tab:discrimination-gap}. Points closer to the origin correspond to
lower error on both measures. The figure is intended as a descriptive complement to
the tables, which report the exact values and Monte Carlo standard errors.}}
\label{fig:flr-mmr-summary}
\end{figure}

\begin{figure}[ht]
\centering
\fbox{
\begin{minipage}{0.90\textwidth}
\centering
Texts
$\longrightarrow$
Embedding [transformer model]
$\longrightarrow$
Scoring [generative/discriminative model]

\vspace{0.8em}

$\longrightarrow$

Separate streams of harder cases and the rest

[e.g.\ by discrimination gap]

\vspace{0.8em}

$\longrightarrow$

\begin{tabular}{ll}
hard cases & $\longrightarrow$ Classification [cross-encoder] \\
rest       & $\longrightarrow$ Classification [generative/discriminative model]
\end{tabular}
\end{minipage}
}
\caption{Workflow of entity resolution combining different models. Type of model in
square brackets.}
\label{fig:combined-workflow}
\end{figure}

Figure \ref{fig:combined-workflow} shows the proposed workflow of entity
resolution with the additional details of how the traditional scoring models can be
combined with the transformer model. In particular, the workflow provides a way to
control the accuracy-cost tradeoff by using the traditional generative or
discriminative models to identify the hard cases for cross-encoder link
classification, while leaving the rest of the link classification to the traditional
model.

The potential computational saving arises because fine-tuning is performed once
for a given labelled sample, whereas the subsequent cost of cross-encoder
classification increases with the number of text pairs sent to it. Once the
bi-encoder embeddings have been obtained, scoring by the traditional models is
comparatively inexpensive. Routing only a subset of cases to the cross-encoder
therefore reduces the number of pairwise cross-encoder evaluations. \rev{As a
hardware-independent cost measure, Table~\ref{tab:combined-procedure} reports the
realised number of cross-encoder pair evaluations for the two training-derived routing
rules relative to sending all $k=50$ candidate pairs to the cross-encoder.} The realised
runtime saving will depend on the routing rule, block size, hardware and
implementation, and is not benchmarked systematically in the present paper.

Although the benefits of creating separate case streams for link classification and the
feasibility of using the traditional models for this task are clear by now, how to determine
exactly which cases should be sent to the cross-encoder model must depend on the cost
constraint and the acceptable accuracy in each application. It is also possible to consider
more refined schemes than, say, choosing a single threshold of the discrimination gap.

For instance, by MEC using a generative model, one can estimate both the FLR and MMR
for the classified link set, and use them to filter out the cases for the cross-encoder model.
Or, to further reduce the cost, one may consider sending fewer candidates for each case
passed on to the cross-encoder model, provided the left-out candidates have much lower
match probabilities than the selected ones according to the scoring model, in which
respect one may consider more refined discrimination measures than the gap
\eqref{eq:discrimination-gap} between the top two candidate match probabilities.

However, we notice that there is another issue related to the exact scheme of creating
separate case streams, which concerns the assessment of linkage errors by transformer
models, as we shall discuss in the next section. 
\rev{To demonstrate the combined
procedure directly, Section~\ref{sec:workable-error-assessment} also reports pooled
FLR and MMR under two routing rules whose discrimination-gap cutoffs are determined
from the training records and then applied unchanged to the test records.}


\subsection{Linkage error assessment of transformer models}
\label{sec:error-assessment}

Assessing the link classification errors of the cross-encoder model faces two major issues.
As we have explained in Section 1.2, the test errors only refer to the model fine-tuned on
the training set, which causes the first issue if one would like to fine-tune the model on the
whole labelled match sample using the hyperparameters settled on the training set, which
is a very common practice in Machine Learning \citep[e.g.][]{hastie2009elements}. The
second issue is related to the creation of separate case streams for link classification in the
proposed workflow, which means, for example, that the cross-encoder model would only be
applied to the selected harder cases instead of all the out-of-sample cases.

Below we investigate the first issue, before proposing a workable approach to linkage error
assessment, even though it does not yield unbiased estimation of the FLR and MMR.


\subsubsection{Sensitivity of training-test split}
\label{sec:split-sensitivity}

We have divided the sampled file-$A$ records into training and test sets 50--50 so far
in this paper. To study sensitivity to this choice, let $A_{s,f}$ denote an SRSWOR sample
containing a fraction $f=0.5,0.6,\ldots,0.9$ of the sampled file-$A$ records $A_s$, and let

\[
s_{M,f}
=
\{(i,j)\in s_M:i\in A_{s,f}\}
\]

denote the labelled match sample induced by those selected records. The cross-encoder is
trained using the records in $A_{s,f}$ and evaluated on the complementary file-$A$ records
$A_s\setminus A_{s,f}$, together with their associated gold-standard matches.

For each value of $f$, we repeat the random split ten times. The cross-encoder is
fine-tuned separately on each training sample using the same model and hyperparameters
as in Section \ref{sec:fine-tuning-models}. Within each training sample, 20\% of the
constructed training pairs are used as validation data to determine the link-classification
threshold by maximising the F1 score. The FLR and MMR are then evaluated on the
corresponding held-out file-$A$ records.

Table \ref{tab:training-fraction} reports the mean error rates over the ten replications,
together with their Monte Carlo standard errors. We also report the relative Monte Carlo
standard error, defined as the Monte Carlo standard error divided by the absolute mean
error rate, as an indication of the precision of the Monte Carlo mean.

\begin{table}[ht]
\centering
\caption{FLR and MMR of the cross-encoder model by training-set fraction $f$, based
on 10 replications. SE denotes the Monte Carlo standard error and RMCSE the relative
Monte Carlo standard error.}
\begin{tabular}{lrrrrrrr}
\toprule
& & \multicolumn{3}{c}{FLR} & \multicolumn{3}{c}{MMR} \\
\cmidrule(lr){3-5}\cmidrule(lr){6-8}
Dataset & $f$ & Mean & SE & RMCSE & Mean & SE & RMCSE \\
\midrule
DBLP-Scholar
& 0.5 & 0.017 & 0.001 & 0.058 & 0.015 & 0.001 & 0.077 \\
& 0.6 & 0.019 & 0.001 & 0.078 & 0.016 & 0.001 & 0.077 \\
& 0.7 & 0.014 & 0.002 & 0.136 & 0.015 & 0.002 & 0.110 \\
& 0.8 & 0.016 & 0.003 & 0.171 & 0.016 & 0.002 & 0.102 \\
& 0.9 & 0.015 & 0.003 & 0.188 & 0.019 & 0.002 & 0.115 \\
\addlinespace
Abt-Buy
& 0.5 & 0.039 & 0.003 & 0.064 & 0.160 & 0.005 & 0.033 \\
& 0.6 & 0.034 & 0.003 & 0.076 & 0.168 & 0.009 & 0.051 \\
& 0.7 & 0.029 & 0.004 & 0.133 & 0.164 & 0.004 & 0.026 \\
& 0.8 & 0.033 & 0.002 & 0.066 & 0.174 & 0.007 & 0.041 \\
& 0.9 & 0.029 & 0.004 & 0.148 & 0.158 & 0.013 & 0.080 \\
\bottomrule
\end{tabular}
\label{tab:training-fraction}
\end{table}

Table \ref{tab:training-fraction} gives the results from ten Monte Carlo
replications of the record-level training-test split for each given $f$. Although the FLR and MMR are not constant with respect
to the training-set size, there is no systematic improvement in either error rate as $f$
increases. For DBLP-Scholar, the mean FLR remains between approximately 0.014 and
0.019 over the five training fractions, while the MMR remains between approximately
0.015 and 0.019. For Abt-Buy, the FLR varies between approximately 0.029 and 0.039,
whereas the MMR varies between approximately 0.158 and 0.174.

The error rates can therefore be considered relatively stable over the range of training
fractions considered here, especially given the variability between random splits. At the
same time, increasing $f$ reduces the size of the holdout test set. For DBLP-Scholar,
the relative Monte Carlo standard error of the FLR becomes noticeably larger at the
higher training fractions. The corresponding RMCSE of the MMR is also somewhat
larger for several of the higher values of $f$, although the pattern is not monotonic.
For Abt-Buy, the relative Monte Carlo standard errors remain comparatively small for
most values of $f$, but increase for both error measures at $f=0.9$. Thus, allocating
more of the labelled sample to fine-tuning does not necessarily yield visibly lower test
error in these experiments, while it leaves fewer records for assessing that error.

However, the stability of the observed error rates does not solve the more fundamental
problem mentioned above. For any $f<1$, the test error in Table \ref{tab:training-fraction}
refers to the cross-encoder fine-tuned using the records in $A_{s,f}$, rather
than to a model subsequently fine-tuned using the complete sampled file-$A$
set $A_s$. Whether
the former can provide a useful estimate of the error of the latter is a separate question,
which we investigate next.


\subsubsection{Bias of subsample test error estimation}
\label{sec:subsample-bias}

We can explore the bias of subsample test error estimation for the larger-sample trained
cross-encoder model by the following setup. Let $A_{s,f}$ and the induced match sample
$s_{M,f}$ be given as above, with $A_s\setminus A_{s,f}$ providing the outer holdout
records. Now, treating $A_{s,f}$ as if it were the whole labelled file-$A$ sample, we create
an inner training sample $A^*_{s,f}\subset A_{s,f}$ and the corresponding inner test records.
Let $s^*_{M,f}$ denote the labelled match sample induced by $A^*_{s,f}$.

The sampling is therefore carried out at the level of file-$A$ records, with their associated
candidate blocks kept together. Let $\mathrm{FLR}$ denote the error of the model trained
using $A_{s,f}$ and evaluated on the outer holdout $A_s\setminus A_{s,f}$, which simulates
the target FLR of interest. Let $\mathrm{FLR}^*$ denote the error of the model trained using
$A^*_{s,f}$ and evaluated on the corresponding inner test records, which simulates
subsample test error estimation in practice. The replicate-level difference

\begin{equation}
\widehat{\mathrm{Bias}}_{\mathrm{FLR}}
=
\mathrm{FLR}^*-\mathrm{FLR}
\end{equation}

can then be used to assess the bias of the subsample test-set estimator. The same
construction is applied to the MMR,

\begin{equation}
\widehat{\mathrm{Bias}}_{\mathrm{MMR}}
=
\mathrm{MMR}^*-\mathrm{MMR}.
\end{equation}

We consider two choices of the inner training records $A^*_{s,f}$ given each
$A_{s,f}$. In addition to splitting $A_{s,f}$ by SRSWOR, we consider bootstrap
resampling, where a sample of size $|A_{s,f}|$ is drawn with replacement from the
file-$A$ records in $A_{s,f}$ and the source records not selected in the bootstrap sample
form the out-of-bag test set. Since a bootstrap sample of size $|A_{s,f}|$ contains on
average $1-e^{-1}\approx0.632$ distinct source records, we use a 65--35 training-test
split when constructing $A^*_{s,f}$ by SRSWOR for comparison.

\begin{table}[ht]
\centering
\caption{Bias of FLR and MMR estimation by sampling fraction $f$, using SRSWOR
or bootstrap subsampling. Results are based on 10 replications. Monte Carlo standard
errors (SE) of the estimated bias are given in parentheses.}
\label{tab:subsample-bias}
\begin{tabular}{llrrr}
\toprule
Dataset & $f$ & Target & SRSWOR bias (SE) & Bootstrap bias (SE) \\
\midrule
\multicolumn{5}{l}{\textit{Panel A: FLR}} \\[2pt]

DBLP-Scholar
& 0.5 & 0.017 & 0.003 (0.003) & 0.004 (0.002) \\
& 0.6 & 0.019 & 0.005 (0.003) & 0.003 (0.002) \\
& 0.7 & 0.014 & 0.005 (0.003) & 0.010 (0.003) \\
& 0.8 & 0.016 & 0.003 (0.004) & 0.008 (0.004) \\
& 0.9 & 0.015 & 0.004 (0.003) & 0.007 (0.003) \\

\addlinespace
Abt-Buy
& 0.5 & 0.039 &  0.001 (0.004) & 0.001 (0.006) \\
& 0.6 & 0.034 & -0.001 (0.006) & 0.004 (0.006) \\
& 0.7 & 0.029 &  0.007 (0.007) & 0.008 (0.006) \\
& 0.8 & 0.033 &  0.001 (0.005) & 0.001 (0.004) \\
& 0.9 & 0.029 &  0.004 (0.005) & 0.006 (0.006) \\

\midrule
\multicolumn{5}{l}{\textit{Panel B: MMR}} \\[2pt]

DBLP-Scholar
& 0.5 & 0.015 & -0.002 (0.001) &  0.005 (0.002) \\
& 0.6 & 0.016 & -0.001 (0.002) &  0.000 (0.002) \\
& 0.7 & 0.015 & -0.001 (0.003) &  0.003 (0.003) \\
& 0.8 & 0.016 &  0.001 (0.003) &  0.002 (0.003) \\
& 0.9 & 0.019 & -0.003 (0.003) & -0.003 (0.002) \\

\addlinespace
Abt-Buy
& 0.5 & 0.160 &  0.015 (0.012) &  0.010 (0.009) \\
& 0.6 & 0.168 &  0.015 (0.014) & -0.004 (0.013) \\
& 0.7 & 0.164 & -0.002 (0.008) &  0.007 (0.009) \\
& 0.8 & 0.174 & -0.007 (0.014) & -0.008 (0.008) \\
& 0.9 & 0.158 &  0.008 (0.018) &  0.007 (0.016) \\

\bottomrule
\end{tabular}
\end{table}

Table \ref{tab:subsample-bias} shows, for each outer training fraction
$f$, the target out-of-sample error and the estimated bias of the corresponding subsample
error under SRSWOR and bootstrap sampling. The results are averaged over ten Monte
Carlo replications, with the Monte Carlo SE of the bias estimator given in parentheses.

There are no large or systematic differences between the SRSWOR and bootstrap results.

\rev{For DBLP-Scholar, the estimated FLR biases are positive for all values of $f$ under both
sampling schemes, but the common sign should not be interpreted as evidence of a systematic
upward bias at this replication count. The estimates range from approximately 0.003 to 0.005
under SRSWOR and from approximately 0.003 to 0.010 under bootstrap sampling, and most
are within about one to two Monte Carlo standard errors of zero. Some of the bootstrap
estimates are larger relative to their Monte Carlo standard errors, but the ten-replication
experiment does not establish a common bias direction. The biases are nevertheless not
negligible relative to the very low target FLR for DBLP-Scholar.}

The corresponding MMR biases for DBLP-Scholar are generally smaller and do not have
a common direction. Under SRSWOR they range from approximately $-0.003$ to 0.001,
while under bootstrap sampling they range from approximately $-0.003$ to 0.005.
Several of the estimates are of similar magnitude to their Monte Carlo standard errors.

The picture is also mixed for Abt-Buy. The FLR biases are generally small and mostly
positive, with the exception of the SRSWOR result at $f=0.6$. The MMR bias changes
sign as $f$ varies under both sampling schemes, and its Monte Carlo standard error is
often large relative to the estimated bias. There is therefore no evidence here of a
common direction of bias that could simply be corrected by adding or subtracting a
fixed amount.

Overall, the results suggest that subsample test errors can provide a rough indication of
the error of a model trained on a larger labelled sample, but they should not be regarded
as unbiased estimates of that error. Nor do the results identify either SRSWOR or
bootstrap sampling as uniformly preferable. The difference between the model trained on
the subsample and the model trained on the larger labelled sample remains part of the
error-assessment problem.

\subsubsection{A workable approach to error assessment}
\label{sec:workable-error-assessment}

In light of the investigation above, a workable approach to linkage error assessment can
be given as follows. Assume the same partition of training and test sets for all the models,
and let $A_{\mathrm{test}}$ denote the file-$A$ records in the test set. Let

\[
A_{\mathrm{test}}^H \cup A_{\mathrm{test}}^R
=
A_{\mathrm{test}},
\qquad
A_{\mathrm{test}}^H \cap A_{\mathrm{test}}^R
=
\varnothing,
\]

denote the separation of hard and rest cases for link classification, where
$A_{\mathrm{test}}^H$ contains those selected for the fine-tuned cross-encoder model and
$A_{\mathrm{test}}^R$ the rest for link classification by the chosen generative or
discriminative model. \rev{For error assessment, the rule producing this partition must
be fixed using the training or validation data before the test set is evaluated; selecting
or tuning the rule from the test-set discrimination gaps or errors would make the
hard/rest split in-sample.} Let

\[
\phi
=
\frac{|A_{\mathrm{test}}^H|}
{|A_{\mathrm{test}}|}
\]

be the test-set proportion selected for cross-encoder classification. This quantity directly
describes the proportion of records routed to the more costly model and is therefore useful
for describing the computational cost of the combined procedure.

To estimate the linkage errors of the combined procedure, however, the component error
rates must be combined according to their respective denominators in Equation (1).
For $g\in\{R,H\}$, let

\[
C_g
=
\sum_{i\in A_{\mathrm{test}}^g\cap A(s_M)}
\ell_i y_{ij_i}
\]

denote the number of correctly linked matchable records in stratum $g$, let

\[
D_g
=
\sum_{i\in A_{\mathrm{test}}^g}\ell_i
\]

denote the number of declared links, and let

\[
M_g
=
|A_{\mathrm{test}}^g\cap A(s_M)|
\]

denote the number of matchable file-$A$ records. The linkage errors of the complete combined procedure are most directly
calculated by pooling the corresponding counts across the two strata. Provided that
at least one link is declared, the FLR and MMR are

\begin{equation}
\label{eq:combined-errors}
\widehat{\mathrm{FLR}}
=
1-\frac{C_R+C_H}{D_R+D_H},
\qquad
\widehat{\mathrm{MMR}}
=
1-\frac{C_R+C_H}{M_R+M_H}.
\end{equation}

This is simply Equation (1) applied to the link decisions produced by the complete
combined procedure.

\rev{To instantiate this procedure, we consider two explicit routing rules. For each
random split, the 25th or 50th percentile of the discrimination gaps among the outer
training records is used as the routing cutoff. The cutoff is therefore fixed before the
test set is evaluated. Test records with a gap at or below the cutoff are sent to the
fine-tuned cross-encoder, while the remaining records are classified by logistic
regression using the fine-tuned bi-encoder similarities. Because the distribution of gaps
may differ between the training and test records, the realised test-set routing proportion
$\phi$ need not equal 0.25 or 0.50. Table~\ref{tab:combined-procedure} reports the
resulting pooled error rates from Equation~\eqref{eq:combined-errors}.}

\begin{table}[ht]
\centering
\caption{\rev{Combined-procedure error rates under training-derived discrimination-gap
routing rules. Entries for $\phi$, FLR and MMR give the mean (Monte Carlo standard
error) over seeds 42--44. Cross-encoder evaluations are the mean number of routed
record pairs per split for $k=50$, relative to evaluating all 50 candidates for every test
record; the percentage in parentheses is the corresponding reduction.}}
\scriptsize
\setlength{\tabcolsep}{4pt}
\begin{tabular}{llrrrr}
\toprule
Dataset & Training-gap rule & $\phi$ & FLR & MMR & Cross-encoder evaluations \\
\midrule
DBLP-Scholar & 25th percentile & \rev{0.246 (0.0033)} & \rev{0.017 (0.0018)} & \rev{0.033 (0.0025)} & \rev{16,117 / 65,400 (75.4\%)} \\
             & 50th percentile & \rev{0.485 (0.0123)} & \rev{0.017 (0.0034)} & \rev{0.024 (0.0035)} & \rev{31,700 / 65,400 (51.5\%)} \\
\addlinespace
Abt-Buy      & 25th percentile & \rev{0.333 (0.0115)} & \rev{0.050 (0.0035)} & \rev{0.148 (0.0048)} & \rev{9,000 / 27,050 (66.7\%)} \\
             & 50th percentile & \rev{0.590 (0.0149)} & \rev{0.042 (0.0021)} & \rev{0.145 (0.0138)} & \rev{15,950 / 27,050 (41.0\%)} \\
\bottomrule
\end{tabular}
\label{tab:combined-procedure}
\end{table}

\rev{The two rules exhibit the intended accuracy-cost tradeoff directly. For
DBLP-Scholar, the 50th-percentile training-gap rule routes 0.485 of the test records on
average and reduces cross-encoder pair evaluations by 51.5\%, with pooled FLR 0.017
and MMR 0.024. The corresponding Abt-Buy rule routes 0.590 of the test records and
reduces cross-encoder evaluations by 41.0\%, with pooled FLR 0.042 and MMR 0.145.
For reference, applying the fine-tuned cross-encoder to all test records in
Table~\ref{tab:fine-tuning} gives FLR/MMR 0.018/0.014 for DBLP-Scholar and
0.041/0.152 for Abt-Buy, whereas fine-tuned logistic regression gives 0.026/0.027 and
0.053/0.269, respectively. The 25th-percentile rule reduces cross-encoder evaluations
further, at the cost of a different balance between FLR and MMR. These results are
illustrative of the two training-derived routing rules rather than an optimisation over
possible cutoffs.}

For a stratum $g\in\{R,H\}$ with positive relevant denominators, the component
error rates can also be written as

\[
\widehat{\mathrm{FLR}}_g
=
1-\frac{C_g}{D_g},
\qquad
\widehat{\mathrm{MMR}}_g
=
1-\frac{C_g}{M_g}.
\]

When the component denominators are positive, Equation
\eqref{eq:combined-errors} is equivalently

\[
\widehat{\mathrm{FLR}}
=
\frac{D_R}{D_R+D_H}
\widehat{\mathrm{FLR}}_{\mathrm{rest}}
+
\frac{D_H}{D_R+D_H}
\widehat{\mathrm{FLR}}_{\mathrm{cross}},
\]

and

\[
\widehat{\mathrm{MMR}}
=
\frac{M_R}{M_R+M_H}
\widehat{\mathrm{MMR}}_{\mathrm{rest}}
+
\frac{M_H}{M_R+M_H}
\widehat{\mathrm{MMR}}_{\mathrm{cross}}.
\]

Writing the pooled counts first also covers the case in which one routed stratum
declares no links, so that its stratum-specific FLR is undefined even though the
FLR of the complete combined procedure remains well-defined.

In other words, treat $A_{\mathrm{test}}^R$ as the auditing sample for the traditional
model and $A_{\mathrm{test}}^H$ as that for the fine-tuned cross-encoder model. The
difference from evaluating either model alone is that the error assessment follows the same
routing rule that would be used in application. Hence the resulting FLR and MMR refer to
the combined linkage procedure actually being considered.

Notice that the routing proportion $\phi$ is not generally the appropriate weight for both
error rates. For the MMR it coincides with the weight in
Equation \eqref{eq:combined-errors} when all test records are matchable. For the FLR,
the relevant weights depend instead on the numbers of links declared by the two models.
This follows directly from the different denominators of FLR and MMR in Equation (1).

The estimates in Equation \eqref{eq:combined-errors} can be used to assess the linkage
errors of the procedure in which the models learned from the training set are subsequently
applied without change outside the labelled sample. This includes the fine-tuned
bi-encoder and cross-encoder models, the chosen generative or discriminative model, and
the rule used to separate the hard cases from the rest.

The matter is less straightforward if some or all of the models are subsequently retrained
using the complete labelled sample before out-of-sample deployment. In that case, the
test-set estimates refer to models trained on less data than the models eventually deployed.
If the full-sample retrained models have lower linkage errors, the test-set estimates may be
conservative. However, the results in Section \ref{sec:subsample-bias} show that the
direction and magnitude of this discrepancy need not be the same for FLR and MMR or
across applications. Bias exploration following the setup of that section can therefore be
used to gather the relevant evidence.

Thus, the approach does not remove the general difficulty of estimating the error of a
model that is retrained after evaluation. It provides a workable error assessment for the
combined entity-resolution procedure while preserving the separation between the
traditional model and the cross-encoder that is required for the accuracy-cost tradeoff.

